\section{GPT-1 y GPT-2: de Fine-Tuning a Zero-Shot Interface}

La evolución temprana de la serie GPT no consiste solo en el aumento de parámetros: también modifica la task interface. GPT-1 aplica supervised fine-tuning después de un preentrenamiento general; GPT-2 pone más énfasis en que la tarea pueda escribirse en un contexto de lenguaje natural para que un mismo language model la resuelva en modo zero-shot.

\subsection{GPT-1: Pretrain then fine-tune}

\term{fine-tuning} (ajuste fino) consiste en seguir actualizando los parámetros preentrenados con datos de la tarea objetivo. GPT-1 usa un Transformer de tipo decoder-style para el generative pretraining y después construye formatos de entrada para tareas como clasificación, entailment y question answering, y ajusta el modelo a cada una.

\lecturefigure{slide-14-gpt1-architecture.jpg}{GPT-1 conecta el preentrenamiento general de lenguaje con el ajuste a cada tarea.}{Video oficial 14:11; Radford et al. 2018.}

\begin{knowledgebox}{Lectura de la figura: reutilizar la arquitectura no significa no modificar nada}
A la izquierda está la Transformer representation unificada; a la derecha, una task-specific input transformation y un output head adaptan el modelo a cada tarea. El pretraining proporciona la inicialización, pero el fine-tuning sigue requiriendo etiquetas, hiperparámetros y un checkpoint por tarea. Es más eficiente que entrenar desde cero, aunque no llega a una prompt-only interface unificada.
\end{knowledgebox}

\subsection{GPT-2: Zero-shot reading comprehension}

Esta sección usa primero reading comprehension para ilustrar la interfaz de texto unificada. \term{zero-shot} significa que, para la tarea objetivo, no se proporcionan examples de entrenamiento ni se actualizan parámetros; solo se pide al modelo que ejecute la tarea mediante el contexto en lenguaje natural. GPT-2 concatena question, context y answer continuation en una secuencia y se observa si el language modeling resuelve implícitamente el QA.

\lecturefigure{slide-15-gpt2-reading-comprehension.jpg}{GPT-2 intenta zero-shot reading comprehension mediante el formato del texto.}{Video oficial 15:21; Radford et al. 2019.}

\begin{knowledgebox}{Lectura de la figura: la tarea se reescribe como continuation}
El modelo no tiene un QA head dedicado; predice el texto que sigue al prompt. La evaluación debe definir answer extraction, exact match o semantic equivalence; si el prompt format cambia mucho, los resultados pueden variar de forma considerable. Zero-shot demuestra la flexibilidad de la interface, pero no implica que el modelo siga realmente cualquier instruction.
\end{knowledgebox}

\subsection{Zero-shot summarization}

La summarization también se escribe como un documento seguido de una indicación y del resumen. El language model aprovecha formatos similares de los datos de entrenamiento y produce una continuation más corta. La dificultad central está en la coverage, la faithfulness y la compression, no en la fluidez gramatical.

\lecturefigure{slide-16-gpt2-summarization.jpg}{GPT-2 replantea la summarization como continuación de texto condicionada.}{Video oficial 16:44.}

\begin{knowledgebox}{Lectura de la figura: la calidad de un resumen tiene al menos tres ejes}
La coverage verifica si se conservan los hechos clave; la faithfulness, si se introducen claims que no están en el texto original; la compression, la longitud y la redundancia. Las overlap metrics como ROUGE solo cubren parte de estas dimensiones. Un modelo generativo puede escribir resúmenes más naturales y, al mismo tiempo, producir con más facilidad errores factuales difíciles de detectar.
\end{knowledgebox}

\subsection{Zero-shot translation}

A continuación, la misma interfaz de continuation se aplica a la traducción. El modelo puede inducirse con «oración fuente + language marker»; los resultados de GPT-2 muestran que una signal multitarea puede absorberse del web text, aunque sigue habiendo una brecha clara frente a los translation systems entrenados específicamente.

\lecturefigure{slide-17-gpt2-translation.jpg}{La zero-shot translation de GPT-2 muestra la signal latente entre idiomas presente en el web pretraining.}{Video oficial 17:30.}

\begin{knowledgebox}{Lectura de la figura: comparar el baseline y la definición de la tarea}
La tendencia de la figura debe contrastarse con un supervised translation baseline. El language model pudo haber visto fragmentos paralelos, plantillas o sitios web similares, por lo que zero-shot no equivale a «nunca haber visto la relación de traducción». Conviene revisar la contamination, la language coverage y los idiomas de pocos recursos, en lugar de mirar solo el promedio de los idiomas con muchos recursos.
\end{knowledgebox}

\begin{warningbox}{Zero-shot es un evaluation protocol, no una propiedad mágica}
Que algo sea zero-shot depende de si la tarea objetivo participó en el entrenamiento, de si el prompt contiene ejemplos, de si el benchmark se filtró y de si se actualizaron los parámetros. Con un web-scale pretraining es difícil demostrar que «la tarea nunca apareció», así que la formulación más precisa es que no hubo supervised fine-tuning explícito sobre el benchmark objetivo.
\end{warningbox}

\subsection{Resumen de la sección}

GPT-1 reutiliza representaciones mediante pretrain/fine-tune; GPT-2 reescribe las tareas como continuation y avanza hacia una interfaz unificada en lenguaje natural. El zero-shot behavior es valioso, pero depende en gran medida del prompt, de la data exposure y de la metric.
